% Part: first-order-logic
% Chapter: models-theories
% Section: set-theory

\documentclass[../../../include/open-logic-section]{subfiles}

\begin{document}

\olfileid[fr]{fol}{mat}{set}

\olsection{La théorie des ensembles}

Presque toutes les mathématiques peuvent être développées dans la
théorie des ensembles. Pour cela, plusieurs conditions doivent être
remplies. Il faut d'abord disposer d'axiomes pour la relation~$\in$.
Différents systèmes d'axiomes ont été élaborés, parfois avec des
propriétés incompatibles de~$\in$. Le système appelé $\Log{ZFC}$,
la théorie des ensembles de Zermelo--Fraenkel avec l'axiome du choix,
se distingue : c'est de loin le plus employé et le plus étudié, car
ses axiomes suffisent à démontrer presque tout ce que les
mathématiciens s'attendent à pouvoir démontrer. Mais, avant de
l'établir, il faut préciser comment nous pouvons même \emph{exprimer}
ce que les mathématiciens souhaitent exprimer. Pour commencer, le
langage ne contient ni !!{constant}s ni !!{function}s. Il n'est donc
pas évident, à première vue, que nous puissions parler d'ensembles
particuliers (comme $\emptyset$ ou $\Nat$), d'opérations sur les
ensembles (comme $X \cup Y$ et $\Pow{X}$), ou encore de constructions
qui font intervenir autre chose que des ensembles, telles que les
relations et les fonctions.

La relation « être un élément de » n'est pas la seule qui nous
intéresse : « être un sous-ensemble de » paraît presque aussi
importante. Nous pouvons toutefois \emph{définir} cette dernière à
partir de la relation d'appartenance. Pour cela, il faut trouver
!!a{formula}~$!A(x, y)$ du langage de la théorie des ensembles,
satisfaite par une paire d'ensembles~$\tuple{X, Y}$ si et seulement
si $X \subseteq Y$. Or $X$ est un sous-ensemble de~$Y$ si et
seulement si tous les !!{element}s de~$X$ sont aussi des
!!{element}s de~$Y$. Nous pouvons donc définir $\subseteq$ par la
formule
\[
\lforall[z][(z \in x \lif z \in y)]
\]
Dès lors, chaque fois que nous voulons employer la relation~$\subseteq$
dans une formule, nous pouvons utiliser à sa place cette formule
(en remplaçant convenablement $x$ et $y$ et, si nécessaire, en
renommant la variable liée~$z$). Par exemple, l'extensionnalité des
ensembles signifie que, si deux ensembles~$x$ et $y$ sont chacun
contenus dans l'autre, alors $x$ et $y$ sont égaux. On peut l'exprimer par
$\lforall[x][\lforall[y][((x \subseteq y \land y
    \subseteq x) \lif x = y)]]$ ou, en substituant la définition
de~$\subseteq$, par
\[
\lforall[x][\lforall[y][((\lforall[z][(z \in x \lif z \in y)] \land
    \lforall[z][(z \in y \lif z \in x)]) \lif x = y)]].
\]
C'est en effet l'un des axiomes de~$\Log{ZFC}$, appelé
« axiome d'extensionnalité ».

Il n'y a pas de !!{constant} pour~$\emptyset$, mais nous pouvons
exprimer « $x$ est vide » par $\lnot \lexists[y][y \in x]$.
« $\emptyset$ existe » devient alors l'!!{sentence}
$\lexists[x][\lnot \lexists[y][y \in x]]$. C'est un autre axiome
de~$\Log{ZFC}$. (Notons que l'axiome d'extensionnalité implique
que l'ensemble vide, s'il existe, est unique.)\footnote{Précision éditoriale : l'extensionnalité seule garantit l'unicité éventuelle, et non l'existence.} Lorsque nous voulons
parler de~$\emptyset$ dans le langage de la théorie des ensembles,
nous écrivons plutôt « il existe un ensemble vide tel que… ».
Par exemple, pour exprimer le fait que $\emptyset$ est un
sous-ensemble de tout ensemble, nous pouvons écrire
\[
\lexists[x][(\lnot\lexists[y][y \in x] \land \lforall[z][x \subseteq
    z])]
\]
où, bien entendu, il faudrait à son tour remplacer $x \subseteq z$
par sa définition.

Pour parler d'opérations sur les ensembles, telles que $X \cup Y$
et $\Pow{X}$, nous procédons de façon semblable. Le langage de la
théorie des ensembles ne comporte aucun symbole de fonction, mais
nous pouvons exprimer les relations fonctionnelles $X \cup Y = Z$
et $\Pow{X} = Y$ par
\begin{align*}
& \lforall[u][((u \in x \lor u \in y) \liff u \in z)]\\
& \lforall[u][(u \subseteq x \liff u \in y )]
\end{align*}
car les !!{element}s de $X \cup Y$ sont exactement les ensembles
qui appartiennent à~$X$ ou à~$Y$, tandis que les !!{element}s de
$\Pow{X}$ sont exactement les sous-ensembles de~$X$. Cela ne nous
permet toutefois pas d'employer $x \cup y$ ou $\Pow{x}$ comme s'il
s'agissait de termes : nous ne pouvons utiliser que les !!{formula}s
entières qui définissent les relations $X \cup Y = Z$ et
$\Pow{X} = Y$. En fait, nous ne savons pas encore si ces relations
sont toujours satisfaites, autrement dit si les réunions et les
ensembles des parties existent toujours. Par exemple, en employant
l'abréviation qui vient d'être définie, l'!!{sentence}
$\lforall[x][\lexists[y][\Pow{x} = y]]$ énonce un autre axiome
de~$\Log{ZFC}$ : l'axiome de l'ensemble des parties.

Qu'en est-il des couples ordonnés et des fonctions ? Nous devons
ici expliquer comment les considérer comme des types particuliers
d'ensembles. Une façon de définir le couple ordonné $\tuple{x, y}$
consiste à poser qu'il s'agit de l'ensemble $\{\{x\}, \{x, y\}\}$.
Mais, comme précédemment, nous ne pouvons introduire !!a{function}
qui désigne cet ensemble; nous ne pouvons définir que la relation
$\tuple{x, y} = z$, c'est-à-dire $\{\{x\}, \{x, y\}\} = z$ :
\[
\lforall[u][(u \in z \liff (\lforall[v][(v \in u \liff v = x)] \lor
  \lforall[v][(v \in u \liff (v = x \lor v = y))]))]
\]
Cette formule dit que les !!{element}s~$u$ de~$z$ sont exactement
les ensembles dont $x$ est l'unique !!{element} ou dont $x$ et~$y$
sont les seuls !!{element}s (autrement dit, les ensembles égaux
à~$\{x\}$ ou à~$\{x, y\}$). Nous pouvons alors exprimer d'autres
relations, par exemple $X \times Y = Z$ :
\[
\lforall[z][(z \in Z \liff \lexists[x][\lexists[y][(x \in
        X \land y \in Y \land \tuple{x, y} = z)]])]
\]

Une fonction $f \colon X \to Y$ peut être conçue comme la relation
$f(x) = y$, c'est-à-dire comme l'ensemble de couples
$\Setabs{\tuple{x,y}}{f(x) = y}$. Nous pouvons alors dire qu'un
ensemble~$f$ est une fonction de~$X$ dans~$Y$ si (a)~c'est une
relation $\subseteq X \times Y$, (b)~elle est totale, c'est-à-dire
que, pour tout $x \in X$, il existe un $y \in Y$ tel que
$\tuple{x, y} \in f$, et (c)~elle est fonctionnelle, c'est-à-dire
que $\tuple{x, y}, \tuple{x, y'} \in f$ implique $y = y'$ (les
valeurs d'une fonction doivent être uniques). Ainsi, « $f$ est une
fonction de~$X$ dans~$Y$ » peut s'écrire :
\begin{align*}
 \lforall[u][(u \in f \lif {}] & \lexists[x][\lexists[y][(x \in X \land y \in
      Y \land \tuple{x, y} = u)]]) \land {}\\
 \lforall[x][(x \in X \lif {}] &
      (\lexists[y][(y \in Y \land \mathrm{maps}(f, x, y))] \land {}\\
& (\lforall[y][\lforall[y'][((\mathrm{maps}(f, x, y) \land
    \mathrm{maps}(f, x, y')) \lif y = y')]])))
\end{align*}
où $\mathrm{maps}(f, x, y)$ abrège $\lexists[v][(v \in f \land
  \tuple{x, y} = v)]$ (cette !!{formula} exprime « $f(x) = y$ »).\footnote{Correction éditoriale : nous fermons la parenthèse extérieure de la condition de totalité et d'unicité, omise dans l'affichage de la source.}

Il n'est pas non plus difficile d'exprimer que $f\colon X \to Y$
est !!{injective}, par exemple :
\begin{multline*}
 f \colon X \to Y \land \lforall[x][\lforall[x'][((x \in X \land x' \in
     X \land {}]] \\
 \lexists[y][(\mathrm{maps}(f, x, y) \land \mathrm{maps}(f,
        x', y))]) \lif x = x')
\end{multline*}
Une fonction~$f\colon X \to Y$ est !!{injective} si et seulement
si, chaque fois que $f$ associe le même~$y$ à deux éléments
$x, x' \in X$, on a $x = x'$. Si nous abrégeons cette formule par
$\mathrm{inj}(f, X, Y)$, nous pouvons déjà énoncer dans le langage
de la théorie des ensembles une proposition aussi peu triviale que
le théorème de Cantor : il n'existe pas de fonction !!{injective}
de $\Pow{X}$ dans~$X$ :
\[
\lforall[X][\lforall[Y][(\Pow{X} = Y \lif
    \lnot\lexists[f][\mathrm{inj}(f, Y, X)])]]
\]

On pourrait croire que la théorie des ensembles a besoin d'un autre
axiome garantissant l'existence d'un ensemble pour chaque propriété
qui définit ses éléments. Si $!A(x)$ est une formule de la théorie
des ensembles dans laquelle la variable~$x$ est libre, considérons
l'!!{sentence}
\[
\lexists[y][\lforall[x][(x \in y \liff !A(x))]].
\]
Elle affirme qu'il existe un ensemble~$y$ dont les !!{element}s
sont tous les $x$ qui satisfont~$!A(x)$, et eux seuls. Ce schéma
s'appelle le « principe de compréhension ». Il paraît fort utile;
malheureusement, il est contradictoire. Prenons
$!A(x) \ident \lnot x \in x$ : le principe de compréhension affirme
alors
\[
\lexists[y][\lforall[x][(x \in y \liff x \notin x)]],
\]
c'est-à-dire l'existence de l'ensemble de tous les ensembles qui ne
sont pas !!{element}s d'eux-mêmes. Aucun tel ensemble ne peut
exister : c'est le paradoxe de Russell. $\Log{ZFC}$ contient en
revanche une version restreinte de ce principe, le schéma de
séparation :
\[
\lforall[z][\lexists[y][\lforall[x][(x \in y \liff (x \in z \land
      !A(x)))]]].
\]
Cette restriction ne conduit pas au paradoxe de Russell.\footnote{Précisions éditoriales : nous rétablissons dans l'affichage la parenthèse extérieure omise dans la source. Dans les schémas de compréhension et de séparation, la variable qui désigne l'ensemble à obtenir ne doit pas apparaître libre dans la formule qui définit ses éléments. Les autres variables libres éventuelles sont des paramètres dont on prend la clôture universelle ; les variables liées auxiliaires sont renommées au besoin.}

\begin{prob}
Montrer que le principe de compréhension est contradictoire en
donnant !!a{derivation} qui établit
\[
\lexists[y][\lforall[x][(x \in y \liff x \notin x)]] \Proves \lfalse.
\]
On pourra d'abord montrer
$(A \lif \lnot A) \land (\lnot A \lif A)
\Proves \lfalse$.
\end{prob}
\end{document}
